\section{Dawn and Full Day: The Shepherds, St.~Anastasia, and the Eternal Word}\label{sec:dawn-day}

Midnight Mass ends in darkness, yet the liturgy of Christmas Day has only begun.
In choral practice, Lauds follows Midnight Mass immediately as light begins to
break, preparing the way for the Second Mass at Dawn and the Third Mass in full
daylight. Across those three movements the Church lifts her gaze from the
shepherds' lantern to the sunrise and finally to the timeless majesty of the
Word before creation.

\subsection{Lauds at Daybreak and the Second Mass at Dawn}

Christmas Lauds opens with a quick, breathless question put to the shepherds
coming back from Bethlehem: \textit{Quem vidistis, pastores? dicite, annuntiate
nobis, in terris quis apparuit?}, ``O ye shepherds, speak, and tell us what ye
have seen; who is appeared in the earth?'' They answer at once: ``We saw the
new-born Child, and angels singing praise to the Lord. Alleluia, Alleluia.'' The
remaining four antiphons marvel at the Virgin Mother who brought forth the King,
the angel's tidings, the heavenly host, and Isaiah's little Child called Mighty
God. Then rises the fifth-century abecedarian hymn of Coelius Sedulius,
\textit{A solis ortus cardine}, ``From where the sunrise had its birth,'' whose
seven stanzas run alphabetically from \textit{A} to \textit{G}, lingering in
stanza~F (\textit{Foeno jacere pertulit}) over the tenderness of Bethlehem: ``On
hay reclined, the Lord Most High, / Within a manger deigned to lie: / And He who
feeds the birds of air / Vouchsafed a little milk to share.''

The Second Mass at Dawn (\textit{Ad secundam Missam in aurora}) keeps its
ancient Roman station \textit{ad Sanctam Anastasiam}, at the foot of the
Palatine Hill (Section~\ref{sec:history}). Its Introit greets the first light of
morning with Isaiah~9:2 and~6: \textit{Lux fulgebit hodie super nos: quia natus
est nobis Dominus}, ``A light shall shine upon us this day: for the Lord is born
to us.'' The Collect asks that the light shining inwardly by faith may break
outward into conduct:

\begin{massbox}{Collect of the Second Mass at Dawn}
\textbf{Latin} (1962 Roman Missal). Da nobis, quaesumus, omnipotens Deus: ut,
qui nova incarnati Verbi tui luce perfundimur; hoc in nostro resplendeat opere,
quod per fidem fulget in mente. Per eundem Dominum nostrum Jesum Christum,
Filium tuum: Qui tecum vivit et regnat in unitate.

\textbf{English} (Lasance, 1945). Grant, we beseech Thee, O almighty God, that
we who are filled with the new light of Thine Incarnate Word, may show forth in
our deeds that which by faith shineth in our minds. Through the same.
\end{massbox}

Immediately after this prayer---at both Low Mass and Sung Mass---the priest adds
the \textbf{Commemoration of St.~Anastasia, Martyr}, whose feast falls on
25~December and whose Palatine basilica originally drew the pope down from
St.~Mary Major at daybreak: ``Grant, we beseech Thee, O almighty God, that we
who honor the solemnity of Thy blessed martyr Anastasia may experience the
effect of her intercession with Thee.'' Even on the birthday of the Lord, the
Roman liturgy refuses to erase a martyr's blood from its memory.

The Epistle (Tit~3:4--7) opens with two words that breathe the entire spirit of
Christmas dawn: ``But when the goodness and kindness of God our Saviour
appeared'' (\textit{Apparuit benignitas et humanitas Salvatoris nostri Dei}),
saving us ``not by the works of justice which we have done, but according to his
mercy, by the laver of regeneration and renovation of the Holy Ghost.'' In his
fifth Homily on Titus, St.~John Chrysostom pauses over that word
``renovation'':

\begin{quote}
For as when a house is in a ruinous state no one places props under it, nor
makes any addition to the old building, but pulls it down to its foundations,
and rebuilds it anew; so in our case, God has not repaired us, but has made us
anew. For this is ``the renewing of the Holy Ghost.'' He has made us new men.
\end{quote}

The Gospel (Luke~2:15--20) follows the shepherds hurrying into Bethlehem to see
``this word that is come to pass,'' finding Mary and Joseph and the infant laid
in the manger, and returning home glorifying God while ``Mary kept all these
words, pondering them in her heart'' (Luke~2:19). In his commentary on Luke,
read as the eighth lesson of Christmas Matins, St.~Ambrose sees the whole
hierarchy of the Church taking shape around that manger: ``Behold the beginning
of the Church. Christ is born, and the shepherds watch; shepherds, to gather
together the scattered sheep of the Gentiles, and to lead them into the fold of
Christ... The flock then is the people, the night is the world, and the
shepherds are the Priests.''

In the Dawn Secret (\textit{Munera nostra}), the priest asks that as he who was
born man shone forth also as God, so this earthly substance may bestow on us
what is divine. Inside the Canon, the proper \textit{Communicantes} now
shifts from \textit{noctem sacratissimam} to \textit{diem sacratissimum},
``keeping that most holy day,'' which will remain at every Mass through
1~January.

\subsection{The Third Mass in the Day and Second Vespers}

When full daylight arrives, the Roman station returns in the 1962 Missal to
\textit{Sancta Maria maior} (inheriting the great papal Mass once sung at
St.~Peter's). The Introit rings out with Isaiah~9:6 and Psalm~97~[98]:1:
\textit{Puer natus est nobis, et filius datus est nobis}, ``A child is born to
us, and a son is given to us, whose government is upon his shoulder: and his
name shall be called, the Angel of great counsel.'' That chant is so closely
wedded to Christmas that the Missal will sing it again on Days~5, 6, and~7
within the Octave, on 1~January, and on the weekdays of 2--5~January.

\begin{massbox}{Collect of the Third Mass in the Day}
\textbf{Latin} (1962 Roman Missal). Concede, quaesumus, omnipotens Deus: ut nos
Unigeniti tui nova per carnem Nativitas liberet; quos sub peccati jugo vetusta
servitus tenet. Per eundem Dominum nostrum Jesum Christum, Filium tuum: Qui
tecum vivit et regnat in unitate Spiritus Sancti, Deus: per omnia saecula
saeculorum. \textbf{R.} Amen.

\textbf{English} (Lasance, 1945). Grant, we beseech Thee, O almighty God, that
the new birth, in the flesh, of Thine only-begotten Son may deliver us whom
slavery from old doth keep under the yoke of sin. Through the same.
\end{massbox}

Every prayer of this Day Mass sets ``new'' against ``old'': the new birth in the
flesh breaking the ancient slavery of sin. The Epistle (Heb~1:1--12) lifts the
eyes straight from Bethlehem to the throne of heaven: God, who spoke of old in
diverse manners by the prophets, ``in these days, hath spoken to us by his Son,
whom he hath appointed heir of all things, by whom also he made the world''---the
brightness of his glory and the figure of his substance, before whom all the
angels of God fall in adoration (Heb~1:1--6). The Gradual and Communion sing
Psalm~97~[98]:3: \textit{Viderunt omnes fines terrae salutare Dei nostri}, ``All
the ends of the earth have seen the salvation of our God,'' while the Alleluia
verse summons the nations: \textit{Dies sanctificatus illuxit nobis: venite,
gentes, et adorate Dominum: quia hodie descendit lux magna super terram}---``A
hallowed day hath dawned for us: come, ye Gentiles, and adore the Lord; for this
day a great light hath descended upon the earth.''

Then comes the sovereign Gospel of Christmas Day, the Prologue of St.~John
(John~1:1--14): ``In the beginning was the Word: and the Word was with God: and
the Word was God... And the Word was made flesh and dwelt among us (and we saw
his glory, the glory as it were of the only begotten of the Father), full of
grace and truth.'' At the words \textit{Et Verbum caro factum est}, priest and
people genuflect together. Reading this Gospel in the ninth lesson of Christmas
Matins (Tractate~1 on John), St.~Augustine warns the believer not to measure
God's Word by fleeting human speech: ``Lest thou shouldest think all things
mean, as thou art accustomed to think of things human, hear and digest
this---`The Word was God.''' Commenting on verse~14 in the passage preserved in
St.~Thomas Aquinas's \textit{Catena Aurea}, St.~John Chrysostom guards the
exact meaning of ``was made flesh'':

\begin{quote}
God's own Son was made the Son of Man, that He might make the sons of men the
sons of God. Now when thou hearest that the Word was made flesh, be not
disturbed, for He did not change His substance into flesh, which it were indeed
impious to suppose; but remaining what He was, took upon Him the form of a
servant. But as there are some who say, that the whole of the incarnation was
only in appearance, to refute such a blasphemy, he used the expression
\textit{was made}, meaning to represent not a conversion of substance, but an
assumption of real flesh.
\end{quote}

The Day Postcommunion (\textit{Praesta, quaesumus}) draws the whole arc of the
three Masses to one point: ``that, as the Saviour of the world, born this day,
is unto us the author of divine generation, so He may also be the bestower of
immortality.''

At dusk on 25~December, Second Vespers gathers the psalms of the Messiah's
eternal kingdom (Psalms~109~[110], 110~[111], 111~[112], 129~[130], and
131~[132]) and crowns Christmas Day with the \textit{Magnificat} antiphon
\textit{Hodie Christus natus est}: ``This day the Christ is born: this day the
Saviour is appeared: this day the Angels sing praise in the earth and the
Archangels rejoice: this day the righteous are glad and say: Glory to God in the
highest. Alleluia.'' Under the 1960 \textit{Codex Rubricarum} (General
Rubrics~108--111), ordinary commemorations never occur at Vespers and
second-class feasts of saints carry no First Vespers; Second Vespers of
Christmas therefore closes in the unbroken radiance of the Nativity (admitting a
single privileged commemoration only of the Sunday within the Octave when
26~December falls on a Sunday). Not until the night office and Lauds of
26~December do the newborn King's closest companions step up to the crib.
